% ─────────────────────────────────────────────────────────────────────────────
%  Chapter 6 — Experimental Evaluation
%  Main benchmark: 10×parallel + 10×serial (all runs retained)
%  Scalability benchmark: 5×N1 + 5×N2 + 5×N4
%  All results fully populated — no placeholders remain.
% ─────────────────────────────────────────────────────────────────────────────

% TikZ styles for this chapter's flowcharts
% (reuse lcXxx from cap5 where possible; add chapter-specific styles here)
\tikzset{
  evstep/.style={rectangle, rounded corners=3pt, draw=black!70, fill=blue!8,
                 text width=5.5cm, minimum height=0.8cm, text centered,
                 font=\small, inner sep=5pt},
  evdec/.style={diamond, aspect=2.5, draw=black!70, fill=orange!15,
                text width=2.8cm, minimum height=0.75cm, text centered,
                font=\small, inner sep=2pt},
  evio/.style={ellipse, draw=black!70, fill=green!20,
               text width=2.8cm, minimum height=0.7cm,
               text centered, font=\small, inner sep=3pt},
  evloop/.style={rectangle, rounded corners=3pt, draw=black!60, fill=gray!8,
                 text width=4.5cm, minimum height=0.75cm, text centered,
                 font=\small, inner sep=5pt},
  evarr/.style={thick,->,>=Stealth},
  evline/.style={thick},
}

% ── Helper macros for placeholders (kept for future chapter extensions) ─────
% \RESULT{mean}{std}{unit}  →  MEAN ± STD UNIT

% ─────────────────────────────────────────────────────────────────────────────
\section{Overview}
\label{sec:eval-overview}
% ─────────────────────────────────────────────────────────────────────────────

This chapter presents the experimental evaluation of the \gls{aerosouda} Remote cluster 
prototype described in Chapter~5. The primary objective is to quantify the performance
benefit of the distributed MapReduce pipeline over a single-executor baseline on the same
Spark infrastructure for the representative production-scale simulation batch used in the
study.

The evaluation is structured as a controlled benchmarking study.  Section~6.1 defines the
experimental setup: the input datasets, analysis scenarios, cluster configurations, and
measurement protocol.  Section~6.2 introduces the performance metrics collected.  The results
for each experimental condition are presented in Sections~6.3 and~6.4, followed by a
discussion of speedup, overhead, and scalability in Section~6.5.

% ─────────────────────────────────────────────────────────────────────────────
\section{Experimental Setup}
\label{sec:eval-setup}
% ─────────────────────────────────────────────────────────────────────────────

\subsection{Input Dataset}
\label{subsec:eval-datasets}

The benchmark uses two simulation batch \gls{csv} datasets:

\begin{itemize}
    \item Quick-Reaction dataset:
          1\,888\,850 rows, approximately 1.3\,GB, involving 10 simulation runs.
          This is the production batch generated from the Quick-Reaction scenario
          described in Section~\ref{sec:cap2-sensitivity-analysis-sa}.
    \item Launcher Longitude dataset:
          2\,290\,288 rows, approximately 1.5\,GB, also involving 10 simulation runs.
          This batch consists in varying the longitude of a \gls{gbad} missile launcher
          across those 10 runs. This dataset's scenario works the same way as the
          \gls{qr} scenario does, only changing the fighters for the \gls{gbad}
          launcher as the opposing force.
\end{itemize}

\subsection{Analysis Scenarios}
\label{subsec:eval-scenarios}

Two \gls{yaml} scenario configurations are evaluated in this chapter:

\begin{enumerate}
    \item Quick-Reaction: computes the Pearson
          correlation between the longitude of the red-side launcher and the survival time
          (maximum frame number) of blue-side fighter entities.  This scenario uses
          \texttt{scope: per\_entity} for the Y variable, producing one data point per
          (entity, experiment) pair.
    \item Launcher Longitude: computes the
          Pearson correlation between the longitude of the red-side \gls{gbad} launcher and
          the maximum damage sustained by blue-side fighter entities.  The
          aggregation field (\texttt{state.damage}) requires deeper nested-dictionary
          traversal than the frame-counting in Scenario~1.
\end{enumerate}

Both scenarios use \texttt{scope: per\_experiment} for the X variable and
\texttt{scope: per\_entity} for the Y variable.  Consequently, one experiment-level X
value can be paired with multiple entity-level Y observations (two blue fighters per
run).  Pearson correlation remains well-defined for these observations, but rows from the
same experiment are not strictly independent; results are therefore interpreted as an
entity-level association measure conditioned on the experiment-level parameter setting.

\subsection{Processing Modes: Serial and Parallel}
\label{subsec:eval-modes}

A central methodological question is how to define a serial baseline that is
comparable to the parallel implementation without requiring a separate codebase.  This work
adopts the following definition:

\begin{itemize}
    \item Serial mode (baseline): the distributed processor runs with a single
          executor, single core, and single \gls{rdd} partition
          (\texttt{executor.instances=1},
          \texttt{executor.cores=1},
          \texttt{--agg-partitions~1}).
          The code path is identical to the parallel mode; only the resource
          configuration differs.
    \item Parallel mode (target): the processor uses four executors, with each
          executor assigned two cores and the \gls{rdd} partitioned into 16
          segments (default configuration).  Four executors execute map tasks concurrently.
\end{itemize}

This baseline definition ensures an apples-to-apples comparison: both modes use the same
Spark infrastructure, the same code, and the same \gls{hdfs} storage layer.  The measured
speedup therefore reflects the benefit of task-level parallelism within the
\gls{aerosouda} architecture, including all fixed costs (\gls{jvm} startup, \gls{hdfs} read, output serialisation).

Using a single Spark executor as the serial baseline also avoids conflating the inherent
overhead of the Spark framework with the overhead attributable to parallelism.  A comparison
against a standalone Python pandas script while informative would introduce confounding
factors such as different memory access patterns and the absence of Spark's partitioned
data model.

\subsection{Measurement Protocol}
\label{subsec:eval-protocol}

Each job execution follows the same protocol, illustrated in the Benchmark Workflow in
Figure~\ref{fig:benchmark-flow}:

\begin{enumerate}
    \item The input \gls{csv} is uploaded to \gls{hdfs} before the timing window begins.
    \item The Spark job is launched remotely via \gls{ssh} by the \texttt{tg-cluster}
          benchmark harness, which executes \texttt{run\_csv\_processor.sh} (or
          \texttt{run\_csv\_processor\_serial.sh}) inside the
          \texttt{spark-master} container.
    \item Two timestamps are recorded: $t_0$ (spark-submit invoked on the
          remote host) and $t_1$ (job process exits on the remote host), giving
          $T_{\text{job}} = t_1 - t_0$.  Additionally, the wall-clock time
          from the local client perspective ($T_{\text{total}}$, including
          \gls{ssh} round-trip overhead) is recorded.  The benchmark harness also
          parses fine-grained phase timing markers emitted by the job
          ($T_{\text{init}}$, $T_{\text{map}}$, $T_{\text{reduce}}$, $T_{\text{output}}$)
          and records peak worker \gls{cpu} and memory via periodic \texttt{docker stats}
          queries; the instrumentation is described in
          Section~\ref{subsec:rm-phase-instrumentation}.
    \item The process is repeated ten times per mode.
          All ten runs per mode are retained for statistical analysis (see
          Section~\ref{subsec:eval-warmup}).
\end{enumerate}

\begin{figure}[htbp]
\centering
\resizebox{0.75\textwidth}{!}{%
\begin{tikzpicture}[x=1cm, y=1cm]

\node[evio] (start) at (0, 0) {Start};

\node[evstep] (setup) at (0,-1.5)
    {Upload CSV to HDFS (one-time)};

\node[evloop, text width=6cm] (formod) at (0,-3.0)
    {for each mode in \{serial, parallel\}};

\node[evloop, text width=6cm] (forrun) at (0,-4.5)
    {for run = 1 \ldots 10};

\node[evstep] (launch) at (0,-6.0)
    {SSH: launch Spark job; record $t_0$};

\node[evstep] (done) at (0,-7.5)
    {Job exits; record $t_1$; compute $T_{\text{job}}, T_{\text{total}}$};

\node[evstep] (record) at (0,-9.0)
    {Append row to benchmark CSV};

\node[evstep, fill=gray!15] (stats) at (0,-10.5)
    {Compute mean \& std over all 10 runs};

\node[evio] (end) at (0,-12.0) {End};

\draw[evarr] (start)  -- (setup);
\draw[evarr] (setup)  -- (formod);
\draw[evarr] (formod) -- (forrun);
\draw[evarr] (forrun) -- (launch);
\draw[evarr] (launch) -- (done);
\draw[evarr] (done)   -- (record);

\draw[evline] (record.east) -- ++(2.5,0) -- ++(0,5.3) -- (forrun.east);
\draw[evline] (stats.east)  -- ++(3.5,0) -- ++(0,8.8) -- (formod.east);

\draw[evarr] (record.south) -- (stats.north);
\draw[evarr] (stats)  -- (end);

\end{tikzpicture}%
}
\caption[Benchmark measurement protocol.]{Benchmark measurement protocol.  The outer loop iterates over the two
         processing modes; the inner loop repeats the job ten times per mode.
         All ten runs per mode are retained after a post-hoc warmup check
         (Section~\ref{subsec:eval-warmup}) found no discardable outlier.}
\label{fig:benchmark-flow}
\end{figure}

% ─────────────────────────────────────────────────────────────────────────────
\section{Performance Metrics}
\label{sec:eval-metrics}
% ─────────────────────────────────────────────────────────────────────────────

Core metrics are collected for each job execution:

\begin{itemize}
    \item Total job time ($T_{\text{job}}$, seconds): wall-clock time measured
          on the remote host from the moment \texttt{spark-submit} is invoked until the
          job process exits.  This metric captures the full end-to-end runtime including
          Spark startup, \gls{hdfs} read, computation, reduce phase, and output write.

    \item Total client time ($T_{\text{total}}$, seconds): wall-clock time as
          seen by the local benchmark client, encompassing $T_{\text{job}}$ plus the
          \gls{ssh} command round-trip overhead.  The difference
          $\Delta_{\text{ssh}} = T_{\text{total}} - T_{\text{job}}$ reflects network
          and shell invocation latency and is expected to be smal overall.

    \item Throughput ($R = N_{\text{rows}} / T_{\text{job}}$, rows\,s$^{-1}$):
          number of simulation rows processed per second.  Allows comparison across
          datasets of different sizes.

    \item Speedup ($S = \bar{T}_{\text{job}}^{\text{serial}} /
          \bar{T}_{\text{job}}^{\text{parallel}}$, dimensionless): ratio of the mean
          serial job time to the mean parallel job time.  The theoretical maximum speedup
          with $N = 4$ executors under Amdahl's law~\parencite{amdahl1967} is
          $S_{\max} = 1 / (f + (1-f)/N)$, where $f$ is the non-parallelisable fraction.

    \item Parallel efficiency ($\eta = S / N$, dimensionless):
          fraction of the theoretical linear speedup achieved in practice.
\end{itemize}

In addition to these core metrics, the benchmark harness records fine-grained phase
timings emitted by the remote Spark job:
\[
T_{\text{init}},\; T_{\text{map}},\; T_{\text{reduce}},\; T_{\text{output}}.
\]
Given Spark's lazy execution model, \(T_{\text{map}}\) includes both \gls{hdfs} read
I/O and distributed payload processing.  A residual term is also computed,
\[
T_{\text{other}} = T_{\text{job}} -
\left(T_{\text{init}} + T_{\text{map}} + T_{\text{reduce}} + T_{\text{output}}\right),
\]
to capture fixed runtime costs not directly attributed to the mapped phases
(for example, job orchestration and scheduler overhead).  Finally, peak worker
\gls{cpu} and memory are sampled via \texttt{docker stats} throughout each run.

% ─────────────────────────────────────────────────────────────────────────────
\section{JVM Warm-up Analysis}
\label{subsec:eval-warmup}
% ─────────────────────────────────────────────────────────────────────────────

Before computing summary statistics, a post-hoc outlier check was performed on
the first run of each mode to test whether \gls{jvm} and \gls{jit} warm-up effects distort the
initial measurements.  Warm-up would manifest as run~1 being noticeably slower than the
remaining runs.

Table~\ref{tab:warmup-check} compares the run-1 job time against the mean of runs~2--10
for each mode.  For the parallel mode, run~1 (76.10\,s) is 0.77\,s above
the subsequent mean (75.33\,s). For the serial mode, run~1 (213.30\,s) is actually
faster by 2.25\,s than the subsequent mean (215.55\,s).

\begin{table}[htbp]
\centering
\caption{\gls{jvm} warm-up check: run~1 vs.\ mean of runs~2--10.}
\label{tab:warmup-check}
\begin{tabular}{lcccc}
\toprule
Mode & Run~1 $T_{\text{job}}$ (s) & Mean runs 2--10 (s) \\
\midrule
Parallel & 76.10 & 75.33 \\
Serial   & 213.30 & 215.55 \\
\bottomrule
\end{tabular}
\end{table}

No warm-up artefact is detected: neither run~1 is a statistical outlier,
and the serial first run is directionally faster, not slower.
Accordingly, all ten runs per mode are retained for computing summary statistics.

% ─────────────────────────────────────────────────────────────────────────────
\section{Results}
\label{sec:eval-results}
% ─────────────────────────────────────────────────────────────────────────────

\subsection{Quick-Reaction Scenario}
\label{subsec:eval-qr}

Executed campaign.
The benchmark was executed using the \texttt{tg-cluster} remote client
(Chapter~5) connected via \gls{ssh} to the AeroSOuDA cluster.
Input: \texttt{results\_quickreaction.csv} (1\,888\,850 rows, $\approx$1.3\,GB) with
scenario \texttt{qr\_correlation.yaml} (Pearson correlation, Quick-Reaction Scenario).
Two modes were measured: parallel (4 executors, default configuration) and
serial (1 executor, 1 core, 1 partition as baseline).  Each mode was repeated
ten times; all ten runs are used for statistics (see Section~\ref{subsec:eval-warmup}).

Tables~\ref{tab:qr-serial-runs} and \ref{tab:qr-parallel-runs} list the per-run
measurements.  Table~\ref{tab:qr-summary} summarises the aggregate statistics.

\begin{table}[htbp]
\centering
\caption{Quick-Reaction, per-run timings (serial mode, 1 executor).}
\label{tab:qr-serial-runs}
\begin{tabular}{cccc}
\toprule
Run & $T_{\text{job}}$ (s) & $T_{\text{total}}$ (s) & $R$ (rows/s) \\
\midrule
1  & 213.30 & 213.66 &  8852.5 \\
2  & 216.29 & 216.63 &  8730.3 \\
3  & 216.80 & 217.15 &  8709.8 \\
4  & 216.78 & 217.15 &  8710.6 \\
5  & 214.26 & 214.67 &  8813.1 \\
6  & 215.72 & 216.06 &  8752.4 \\
7  & 215.78 & 216.12 &  8750.0 \\
8  & 215.79 & 216.12 &  8749.9 \\
9  & 213.78 & 214.16 &  8832.6 \\
10 & 214.78 & 218.37 &  8792.1 \\
\bottomrule
\end{tabular}
\end{table}

\begin{table}[htbp]
\centering
\caption{Quick-Reaction, per-run timings (parallel mode, 4 executors).}
\label{tab:qr-parallel-runs}
\begin{tabular}{cccc}
\toprule
Run & $T_{\text{job}}$ (s) & $T_{\text{total}}$ (s) & $R$ (rows/s) \\
\midrule
1  & 76.10 & 76.48 & 24820.6 \\
2  & 74.62 & 74.94 & 25312.9 \\
3  & 76.10 & 76.44 & 24820.6 \\
4  & 75.10 & 75.50 & 25151.1 \\
5  & 75.60 & 75.95 & 24984.8 \\
6  & 74.59 & 74.93 & 25323.1 \\
7  & 75.60 & 75.94 & 24984.8 \\
8  & 75.61 & 75.98 & 24981.5 \\
9  & 75.61 & 75.94 & 24981.5 \\
10 & 75.17 & 75.51 & 25128.3 \\
\bottomrule
\end{tabular}
\end{table}

\begin{table}[htbp]
\centering
\caption{Quick-Reaction, summary statistics (all 10 runs per mode).}
\label{tab:qr-summary}
\begin{tabular}{lcc}
\toprule
Metric & Serial (mean $\pm$ std) & Parallel (mean $\pm$ std) \\
\midrule
$T_{\text{job}}$ (s)    & $215.33 \pm 1.24$   & $75.41 \pm 0.53$    \\
$T_{\text{total}}$ (s)  & $216.01 \pm 1.47$   & $75.76 \pm 0.54$    \\
$\Delta_{\text{ssh}}$ (s) & $0.68 \pm 1.02$   & $0.35 \pm 0.03$     \\
$R$ (rows\,s$^{-1}$)    & $8\,769 \pm 50$   & $25\,049 \pm 177$     \\
\midrule
\multicolumn{3}{l}{Derived metrics} \\
Speedup $S$             & \multicolumn{2}{c}{$2.855$} \\
Parallel efficiency $\eta$ & \multicolumn{2}{c}{$0.714$ (71.4\%)} \\
\bottomrule
\end{tabular}
\end{table}

% ─────────────────────────────────────────────────────────────────────────────
\subsection{Launcher Longitude Scenario}
\label{subsec:eval-ll}
% ─────────────────────────────────────────────────────────────────────────────

Executed campaign.
The Launcher Longitude benchmark was executed using the standard main-benchmark
protocol: ten serial runs (1 executor) and ten parallel runs (4 executors).
Input: \texttt{results\_gbad\_launcher\_longitude.csv} (2\,290\,288 rows, $\approx$1.5\,GB) with
scenario \texttt{analyze\_fighter\_damage\_by\_launcher\_longitude.yaml} (Pearson correlation
between \gls{gbad} launcher longitude and the damage taken by the fighter jets).
Both modes were repeated ten times; all ten runs are used for statistics.

Timing results.
Tables~\ref{tab:ll-serial-runs} and \ref{tab:ll-parallel-runs} list the per-run
measurements.  Table~\ref{tab:ll-summary} summarises the aggregate statistics.

\begin{table}[htbp]
\centering
\caption{Launcher Longitude, per-run timings (serial mode, 1 executor).}
\label{tab:ll-serial-runs}
\begin{tabular}{cccc}
\toprule
Run & $T_{\text{job}}$ (s) & $T_{\text{total}}$ (s) & $R$ (rows/s) \\
\midrule
1  & 250.81 & 251.65 & 9\,132 \\
2  & 252.81 & 253.58 & 9\,059 \\
3  & 253.82 & 254.64 & 9\,023 \\
4  & 251.81 & 252.65 & 9\,095 \\
5  & 250.30 & 251.08 & 9\,150 \\
6  & 252.26 & 253.08 & 9\,079 \\
7  & 252.83 & 253.69 & 9\,059 \\
8  & 252.29 & 253.09 & 9\,078 \\
9  & 250.23 & 251.05 & 9\,153 \\
10 & 253.24 & 254.13 & 9\,044 \\
\bottomrule
\end{tabular}
\end{table}

\begin{table}[htbp]
\centering
\caption{Launcher Longitude, per-run timings (parallel mode, 4 executors).}
\label{tab:ll-parallel-runs}
\begin{tabular}{cccc}
\toprule
Run & $T_{\text{job}}$ (s) & $T_{\text{total}}$ (s) & $R$ (rows/s) \\
\midrule
1  & 81.61 & 82.45 & 28\,064 \\
2  & 79.60 & 80.47 & 28\,772 \\
3  & 80.11 & 80.92 & 28\,589 \\
4  & 79.59 & 80.42 & 28\,776 \\
5  & 80.10 & 81.01 & 28\,593 \\
6  & 79.61 & 80.44 & 28\,769 \\
7  & 81.11 & 80.96 & 28\,237 \\
8  & 80.11 & 80.96 & 28\,589 \\
9  & 79.61 & 80.47 & 28\,769 \\
10 & 80.11 & 80.91 & 28\,589 \\
\bottomrule
\end{tabular}
\end{table}

\begin{table}[htbp]
\centering
\caption{Launcher Longitude, summary statistics (all 10 runs per mode).}
\label{tab:ll-summary}
\begin{tabular}{lcc}
\toprule
Metric & Serial (mean $\pm$ std) & Parallel (mean $\pm$ std) \\
\midrule
$T_{\text{job}}$ (s)      & $252.04 \pm 1.24$   & $80.16 \pm 0.69$    \\
$T_{\text{total}}$ (s)    & $252.86 \pm 1.21$   & $80.90 \pm 0.57$    \\
$\Delta_{\text{ssh}}$ (s) & $0.82 \pm 0.04$     & $0.75 \pm 0.32$     \\
$R$ (rows\,s$^{-1}$)      & $9\,087 \pm 42$     & $28\,573 \pm 243$   \\
\midrule
\multicolumn{3}{l}{Derived metrics} \\
Speedup $S$                & \multicolumn{2}{c}{$3.144$} \\
Parallel efficiency $\eta$ & \multicolumn{2}{c}{$0.786$ (78.6\%)} \\
\bottomrule
\end{tabular}
\end{table}

The Launcher Longitude scenario processes a larger dataset than Quick-Reaction
(2.29M vs 1.89M rows) and achieves a higher speedup ($3.144$ vs $2.855$) and
parallel efficiency ($78.6\%$ vs $71.4\%$).  This is consistent with Amdahl-type
behaviour: a larger parallelisable workload reduces the relative weight of the
non-parallelisable serial component (Spark session startup, driver initialisation, output
serialisation), thereby yielding a higher observed speedup.
Section~\ref{sec:eval-scalability} characterises this behaviour in detail using a
scalability experiment over the same input file.

% ─────────────────────────────────────────────────────────────────────────────
\section{Discussion}
\label{sec:eval-discussion}
% ─────────────────────────────────────────────────────────────────────────────

\subsection{Speedup Analysis}
\label{subsec:eval-speedup}

Table~\ref{tab:speedup-comparison} summarises the speedup and efficiency achieved by
each campaign.

\begin{table}[htbp]
\centering
\caption{Speedup and efficiency comparison across the two benchmark campaigns.}
\label{tab:speedup-comparison}
\begin{tabular}{lcccc}
\toprule
Campaign & Dataset rows & $\bar{T}_{\text{serial}}$ (s) &
$S$ & $\eta$ \\
\midrule
Quick-Reaction   & 1\,888\,850 & 215.33 & 2.855 & 71.4\% \\
Launcher Longitude & 2\,290\,288 & 252.04 & 3.144 & 78.6\% \\
\bottomrule
\end{tabular}
\end{table}

Both campaigns achieve meaningful parallel acceleration on a single-host Docker
deployment.  The higher speedup and efficiency of the Launcher Longitude campaign
($S=3.144$, $\eta=78.6\%$) relative to Quick-Reaction ($S=2.855$, $\eta=71.4\%$)
is consistent with the larger input volume: the additional $\approx400\,000$ rows
increase the parallelisable fraction of the job, reducing the relative overhead of
the fixed serial components and allowing the four executors to deliver a higher net
acceleration.  The key take-away is that the \gls{aerosouda} distributed pipeline
reduces wall-clock time by $65$--$75\%$ relative to the single-executor baseline,
independently of which analytical scenario or dataset is used.

\subsection{Amdahl's Law and Overhead Analysis}
\label{subsec:eval-amdahl}

Amdahl's law~\parencite{amdahl1967} predicts that for a programme with non-parallelisable
fraction $f$, the speedup with $N$ executors is
$S(N)=1/(f+(1-f)/N)$.

In this work, \(f\) is estimated empirically from the Launcher Longitude
scalability timings (Section~\ref{sec:eval-scalability}) by fitting \(T(N)=T_s+T_p/N\),
where \(T_s\) is the serial component and \(T_p\) is the parallelisable component.
Using \(N\in\{1,2,4\}\), the fit yields \(\hat{f}=T_s/(T_s+T_p)\approx0.077\), which
corresponds to \(S_{\max}(2)\approx1.857\) and \(S_{\max}(4)\approx3.247\).  These values
are close to the measured speedups (\(1.867\) and \(3.232\), respectively), confirming
that the observed scaling is well-described by Amdahl-type behaviour under shared-host
contention.

\subsection{Throughput Scaling}
\label{subsec:eval-throughput}

Table~\ref{tab:throughput-summary} summarises throughput for both campaigns.  For
Quick-Reaction, throughput increased from $8\,769 \pm 50$\,rows/s to
$25\,049 \pm 177$\,rows/s, a factor of $2.86\times$.  For Launcher Longitude, it
increased from $9\,087 \pm 42$\,rows/s to $28\,573 \pm 243$\,rows/s, a factor of
$3.14\times$, matching the respective speedup values.

\begin{table}[htbp]
\centering
\caption{Throughput comparison across both benchmark campaigns (all 10 runs per mode).}
\label{tab:throughput-summary}
\begin{tabular}{llccc}
\toprule
Campaign & Mode & Mean (rows/s) & Std (rows/s) \\
\midrule
Quick-Reaction   & Serial   & 8\,769  & 50  \\
                 & Parallel & 25\,049 & 177 \\
\addlinespace
Launcher Longitude & Serial   & 9\,087  & 42  \\
                   & Parallel & 28\,573 & 243 \\
\bottomrule
\end{tabular}
\end{table}

\subsection{SSH Communication Overhead}
\label{subsec:eval-ssh}

The difference $\Delta_{\text{ssh}} = T_{\text{total}} - T_{\text{job}}$ captures the
round-trip cost of the \gls{ssh} command dispatch used by the \texttt{tg-cluster} client.
For the Quick-Reaction parallel mode, the mean overhead is $0.35 \pm 0.03$\,s  
less than $0.5\%$ of total job time.  For the Launcher Longitude parallel mode, the mean
overhead is $0.75 \pm 0.32$\,s, less than $1\%$ of total job time.  In both cases,
\gls{ssh} dispatch adds negligible latency for the distributed execution model.

\subsection{Threats to Validity}
\label{subsec:eval-threats}

A couple of factors limit the generalisability of these results:

\begin{itemize}
    \item Single-host deployment.  All containers (master, four workers,
          namenode, datanodes, history server) ran on the same physical machine.
          The \gls{hdfs} data-locality benefit, a primary architectural advantage  
          cannot be realised in this configuration.  On a genuine multi-node cluster,
          network-aware data placement would reduce inter-node I/O and is expected to
          improve speedup further.
    \item Shared system resources.  Worker containers compete for the same
          physical CPU cores, memory, and disk I/O, introducing contention that would
          not exist on dedicated hardware and limits efficiency at higher executor counts.
          The phase-timing instrumentation described in
          Section~\ref{subsec:rm-phase-instrumentation} enables future benchmark runs to
          quantify each phase's contribution to $T_{\text{job}}$ and to correlate
          worker \gls{cpu} saturation (measured via \texttt{docker stats}) with the
          observed efficiency drop from $\eta=93.4\%$ at $N=2$ to $\eta=80.8\%$ at $N=4$,
          providing an empirical basis for the shared-resource hypothesis.
    \item Dataset coverage.  Two datasets were evaluated: the Quick-Reaction
          batch (1.3\,GB, 1.89M rows) and the Launcher Longitude batch (1.5\,GB, 2.29M
          rows), both consisting in only 10 simulations due to input gathering limitations.
          Additional campaigns with broader data volumes and different entity schemas would
          strongly help generalizing those results.
\end{itemize}

% ─────────────────────────────────────────────────────────────────────────────
\section{Scalability Analysis}
\label{sec:eval-scalability}
% ─────────────────────────────────────────────────────────────────────────────

\subsection{Motivation}
\label{subsec:eval-scalability-motivation}

The main benchmark of Section~\ref{sec:eval-results} provides two data points on
the speedup curve: N=1 (serial baseline) and N=4 (parallel mode).  A single pair of
points is sufficient to claim that parallelism is beneficial, but it does not reveal
whether the system scales efficiently as the number of executors grows.  This section
closes that gap by running the Launcher Longitude full-dataset campaign with N=1, 2, and 4
executors and comparing the resulting speedup curve against the Amdahl's law prediction.

\subsection{Experimental Setup}
\label{subsec:eval-scalability-setup}

The scalability experiment uses the Launcher Longitude input file
(\texttt{results\_gbad\_launcher\_longitude.csv}, \texttt{analyze\_fighter\_damage\_by\_launcher\_longitude.yaml}).
Three executor configurations are evaluated:

\begin{itemize}
    \item N=1: 1 executor, 1 core, 1 partition.
    \item N=2: 2 executors, 2 cores maximum.
    \item N=4: 4 executors, 4 cores maximum.
\end{itemize}

Each configuration is executed with the \texttt{tg-cluster} benchmark harness using
the \texttt{--mode scalability} flag:

\begin{verbatim}
python benchmark.py --runs 5 --mode scalability \
    --csv results_gbad_launcher_longitude.csv --no-download
\end{verbatim}

Five runs per configuration are used.  One run at N=2 produced an execution time of
$882.65$\,s, an outlier attributable to a transient cluster scheduling hiccup  
and was excluded; the remaining four N=2 runs were used for statistics.

\subsection{Theoretical Prediction}
\label{subsec:eval-scalability-theory}

Using the fitted non-parallelisable fraction $\hat{f} \approx 0.077$ estimated from the
scalability timings (Section~\ref{subsec:eval-amdahl}), Amdahl's law predicts:
\[
S_{\max}(N) = \frac{1}{\hat{f} + (1-\hat{f})/N}
\approx \frac{1}{0.077 + 0.923/N}.
\]

\begin{table}[htbp]
\centering
\caption{Amdahl's law predictions vs.\ measured speedup for the scalability experiment
         (Launcher Longitude dataset).}
\label{tab:scalability-amdahl}
\begin{tabular}{cccc}
\toprule
N & $S_{\max}$ (Amdahl) & $S$ (measured) & $\eta$ (measured) \\
\midrule
1 & 1.000 & 1.000 & — \\
2 & 1.857 & 1.867 & 93.4\% \\
4 & 3.247 & 3.232 & 80.8\% \\
\bottomrule
\end{tabular}
\end{table}

\subsection{Results}
\label{subsec:eval-scalability-results}

\begin{table}[htbp]
\centering
\caption{Scalability experiment, mean $T_{\text{job}}$ and speedup by executor count
         (5 runs per configuration, Launcher Longitude dataset).}
\label{tab:scalability-results}
\begin{tabular}{ccccc}
\toprule
N & $\bar{T}_{\text{job}}$ (s) & $\sigma$ (s) & Speedup $S$ & Efficiency $\eta$ \\
\midrule
1 & 260.55 & 4.39 & 1.000 & — \\
2 & 139.55 & 0.75 & 1.867 & 93.4\% \\
4 & 80.61  & 0.87 & 3.232 & 80.8\% \\
\bottomrule
\end{tabular}
\end{table}

\noindent The N=4 result is consistent with the 10-run main benchmark
(Table~\ref{tab:ll-summary}), confirming measurement stability across independent
sessions.  The N=1 mean ($260.55$\,s) is slightly higher than the main benchmark
serial mean ($252.04$\,s), which is expected as the scalability campaign ran under
different system-load conditions.

\subsection{Discussion}
\label{subsec:eval-scalability-discussion}

The results in Table~\ref{tab:scalability-amdahl} exhibit a clear and expected trend:
speedup increases with $N$ but at a diminishing rate, consistent with Amdahl's law.
Two observations are particularly relevant for this thesis.

High efficiency at N=2.
With two executors, the system achieves $S = 1.867$ and $\eta = 93.4\%$, very close to
the fitted Amdahl prediction ($S_{\max}(2)=1.857$).  At this scale, the shared-host
deployment introduces limited contention and the job is nearly perfectly parallel.

Efficiency drop from N=2 to N=4.
The efficiency falls from $93.4\%$ to $80.8\%$ as the executor count doubles.  This
degradation is consistent with shared-host contention: all four worker containers
compete for the same physical cores, memory bus, and disk I/O.  Even so, quadrupling
the executor count reduces wall-clock time by a factor of $3.23\times$, achieving the
highest speedup measured in this thesis.

Cross-dataset consistency.
The Quick-Reaction campaign (Section~\ref{subsec:eval-qr}) uses a smaller dataset
(1.89M rows) and shows a measured speedup of $2.855$ at N=4 with $71.4\%$ efficiency.
The Launcher Longitude campaign uses a larger dataset (2.29M rows) and achieves
$3.144$ at N=4 from the main benchmark and $3.232$ from the scalability experiment.
The increase in speedup with dataset size is consistent with Amdahl's law: a larger
parallelisable workload reduces the relative serial overhead fraction, raising the
efficiency ceiling.  Taken together, the two campaigns confirm that the distributed
pipeline scales beneficially across different input volumes and analytical configurations.
